\section{Evaluation}


\subsection{Agent Evaluation}\label{eval1}


In order to assess the performance of CAMEL (Cooperative Role-playing Communication), we conduct two types of evaluations: (1) Human evaluation, and (2) GPT4 evaluation. We randomly select 100 tasks from our AI Society dataset for evaluation and 100 tasks from our Code dataset. Then, we employ the GPT4 model to summarize the content of the CAMEL conversation-based solution, presenting a consolidated final solution. Particularly, a GPT4 is used since it possesses a larger token limit which is suitable for summarization. Summarization also makes CAMEL agents' solution undetectable by its format, allowing for a more fair comparison. Subsequently, this solution is compared with a single-shot solution generated by the \texttt{gpt-3.5-turbo} model for the same task. Sample tasks are provided in the Appendix.


\textbf{Human Evaluation.} For this evaluation, we present both the CAMEL summarized agent solution and the \texttt{gpt-3.5-turbo} single-shot solution side-by-side to human participants. The identity behind each solution is not revealed. Participants are then asked to vote on whether one solution is superior to the other or if they are equally good. A total of 453 responses were collected during this evaluation. Note that, human evaluation is only done for AI Society, as assessing code is generally harder for humans (without running the code).

\textbf{GPT4 Evaluation.} We engage a GPT4 agent to evaluate the effectiveness of Model 1 (CAMEL Agent solution) versus Model 2 (\texttt{gpt-3.5-turbo} single-shot solution) for each task. More specifically, we prompt GPT4 to score and decide which solution of the two solutions is better.

\textbf{Results.} The summarized results of each evaluation are outlined in Table \ref{agent_eval} which showcases that the CAMEL solution outperforms \texttt{gpt-3.5-turbo} single-shot solution in both the human evaluation and the GPT4 evaluation by a big margin. It is also worth noting that both human evaluation and GPT4 evaluation are highly aligned. 


\begin{table}[h]
  \centering
  \caption{\textbf{Agent Evaluation Results}: Results of the evaluations of the CAMEL agent against \texttt{gpt-3.5-turbo} using both human evaluators and GPT4 consistently show that utilizing a multi-agent cooperative approach is more effective than \texttt{gpt-3.5-turbo}'s single shot solution.}
  \label{agent_eval}
\resizebox{0.75\columnwidth}{!}{%
\begin{tabular}{@{}ccccc@{}}
\toprule
\textbf{Dataset}                     & \textbf{Evaluation Type}  & \textbf{Draw} & \textit{\textbf{gpt-3.5-turbo Wins}} & \textbf{CAMEL Agents Win} \\ \midrule
\multirow{2}{*}{\textbf{AI Society}} & \textbf{Human Evaluation} & 13.3\%        & 10.4\%                               & \textbf{76.3}\%                    \\
                                     & \textbf{GPT4 Evaluation}  & 4.0\%         & 23.0\%                               & \textbf{73.0}\%                    \\ \midrule
\textbf{Code}                        & \textbf{GPT4 Evaluation}  & 0.0\%         & 24.0\%                               & \textbf{76.0}\%                    \\ \bottomrule
\end{tabular}
  }

\end{table}


\subsection{GPT4 for ChatBot Evaluation}\label{sec:emergence}

In this section, we progressively fine-tune a LLaMA 7B model on our generated datasets. By progressively incorporating diverse datasets like AI society, code, math, and science, we expect fine-tuned model to demonstrate the ability to develop an increasingly sophisticated understanding of these domains.

We initially start by training on AI society dataset, which aims to let the model learn about human interactions and societal dynamics. As additional datasets were introduced, such as code, the model gained knowledge of programming logic and syntax, enabling it to generate coherent and executable code snippets. The inclusion of the math dataset further expanded the model's capabilities, allowing it to solve complex equations, reason about abstract concepts, and perform precise calculations. Finally, exposure to the science dataset broadened the model's understanding of scientific theories, empirical observations, and experimental methods. The emergence of model capabilities is measured by evaluating the quality of the model responses, before and after training on the new domain, on a set of questions of varying difficulties from each domain. More precisely, the model is tested on 20 AI Society related tasks, 20 coding tasks, 20 math tasks and 60 science tasks.

Those results are highlighted in Table \ref{emergence} where we see that each time we add a dataset, the model performs better on the incorporated domain. Note that to measure the quality of the models' responses, we follow the evaluation from Section \ref{eval1}, which involves prompting a GPT4 agent to score and decide which solution is better. It is worth noting that an improvement on other domains is also observed in some cases such as when we train on Code we improve on Science. This is because our Code dataset contains problems that solve tasks in particular domains which include scientific domain. Similarly, training on AI Society improves code as AI Society contains the role of a "programmer" and hence coding related conversations. Finally, note that the draws observed in LLaMA-7B vs AI Society in Math reflects equally bad solutions compared to the draws observed in AI Society + Code + Math vs AI Society + Code + Math + Science where the draws are equally good solutions. This progression from AI society to code to math to science highlights the potential of AI models to acquire a versatile and adaptable knowledge base, paralleling the way humans gain expertise in diverse subjects. Sample tasks are provided in the Appendix.

\begin{table}[h!]
\vspace{-0.4cm}
\caption{\textbf{Emergence of Knowledge.} By progressively fine-tuning LLaMA on datasets from different domains, we observe the emergence of knowledge as the model transitions from AI society to code, math, and science. This finding is indicated by the fact that Model 2 almost always performs better than Model 1, especially on the added dataset.}
\label{emergence}
\centering
\resizebox{0.9\columnwidth}{!}{%
\begin{tabular}{c|cccc|cccc|ccc}
\hline
\multirow{2}{*}{\textbf{Dataset}} & \multicolumn{4}{c|}{\textbf{Model 1}}       & \multicolumn{4}{c|}{\textbf{Model 2}}       & \multicolumn{1}{c}{\multirow{2}{*}{\textbf{Draw}}} & \multicolumn{1}{c}{\multirow{2}{*}{\textbf{Model 1}}} & \multicolumn{1}{c}{\multirow{2}{*}{\textbf{Model 2}}} \\
                         & \textbf{AI Society} & \textbf{Code} & \textbf{Math} & \textbf{Science} & \textbf{AI Society} & \textbf{Code} & \textbf{Math} & \textbf{Science} & \multicolumn{1}{c}{}                      & \multicolumn{1}{c}{}                              & \multicolumn{1}{c}{}                              \\ \hline 
\textbf{AI Society}               &            &      &    &         & \cmark     &      &    &         &        0 & 6 & \textbf{14} \\
\textbf{Code}                     &            &      &    &         & \cmark          &      &      &         &   0 & 0 & \textbf{20} \\
\textbf{Math}                     &            &      &    &         & \cmark           &     &      &         &  \textbf{9} & 5 & 6 \\
\textbf{Science}                  &            &      &    &         &  \cmark          &     &      &         &  0 & 13 & \textbf{47} \\ \hline
\textbf{AI Society}               & \cmark     &      &    &         &     \cmark      &  \cmark &  & & 4 & 8 & \textbf{8} \\
\textbf{Code}                     & \cmark     &      &    &         &       \cmark    &  \cmark &  & & 1 & 9 & \textbf{10} \\
\textbf{Math}                     & \cmark     &      &    &         &       \cmark    &  \cmark    & & & 5 & \textbf{8}  & 7 \\
\textbf{Science}                  & \cmark     &      &    &         &      \cmark     &  \cmark   & & & 1 & 19 & \textbf{40} \\ \hline
\textbf{AI Society}               & \cmark     &   \cmark  &         &         &  \cmark   &   \cmark    &   \cmark  &  & 5 & 6 & \textbf{9}     \\
\textbf{Code}                     & \cmark     &   \cmark  &         &         &  \cmark    &  \cmark    &   \cmark  &  & 1 & 9 & \textbf{10} \\
\textbf{Math}                     & \cmark     &   \cmark  &         &         &  \cmark    &  \cmark    &   \cmark  &  & 1 & 3 & \textbf{16}\\
\textbf{Science}                  & \cmark     &   \cmark &      &         &   \cmark         &  \cmark    & \cmark  &  & 3 & 8 & \textbf{49} \\ \hline
\textbf{AI Society}               & \cmark     &   \cmark &   \cmark   &         &  \cmark   &   \cmark    &   \cmark  & \cmark & 3 & 1 & \textbf{16}     \\
\textbf{Code}                     & \cmark     &   \cmark &   \cmark   &         &  \cmark    &  \cmark    &   \cmark  & \cmark & 1 & 8 & \textbf{11}\\
\textbf{Math}                     & \cmark     &   \cmark &   \cmark  &         &  \cmark    &  \cmark    &   \cmark   & \cmark & \textbf{10} & 5 & 5\\
\textbf{Science}                  & \cmark     &   \cmark &   \cmark &         &   \cmark         &  \cmark    &   \cmark   & \cmark & 9 & 2 & \textbf{49}\\ \hline
\textbf{AI Society}               &            &    &      &         &  \cmark   &   \cmark    &   \cmark  & \cmark & 0 & 0 & \textbf{20}     \\
\textbf{Code}                     &            &    &      &         &  \cmark    &  \cmark    &   \cmark   & \cmark & 0 & 0 & \textbf{20} \\
\textbf{Math}                     &            &    &      &         &  \cmark    &  \cmark    &   \cmark   & \cmark & 0 & 0 & \textbf{20}\\
\textbf{Science}                  &            &    &      &         &   \cmark         &  \cmark    &   \cmark   & \cmark & 0 & 0 & \textbf{60}\\ \hline
\end{tabular}}
\end{table}

\subsection{HumanEval$^{(+)}$}
\begin{table}[h]
  \centering
  \caption{\textbf{HumanEval$^{(+)}$ for Various Models.} We test our CAMEL model, which is a LLaMa-7B fine-tuned on all our datasets (AI Society, Code, Math, Science) on HumanEval and HumanEval$^{+}$ benchmarks, where we show competitive pass$@k$ scores with LLaMa-7B and Vicuna-7B.}
  \label{humaneval}
  \scalebox{0.85}{
    \begin{tabular}{ccccc}
      \cline{2-5}
      \multicolumn{1}{l}{} & \multicolumn{2}{l}{HumanEval}  & \multicolumn{2}{l}{HumanEval$^{+}$} \\ \hline
      \textbf{pass}$@k$ [\%]   & $k=1$  & $k=100$ & $k=1$    & $k=100$    \\ \hline
      \textcolor{gray}{\texttt{gpt-3.5-turbo}} & \textcolor{gray}{$69.4$} & \textcolor{gray}{$94.0$} & \textcolor{gray}{$61.7$} & \textcolor{gray}{$89.8$}  \\ \hline
      \textbf{LLaMA-7B}       & $10.5$          & $36.5$           & -               & -                 \\
      \textbf{Vicuna-7B}      & $11.0$          & $42.9$           & $9.9$    & $34.7$              \\
      \textbf{CAMEL-7B}       & $\boldsymbol{14.0}$ & $\boldsymbol{57.9}$  & $\boldsymbol{12.2}$             & $\boldsymbol{50.0}$     \\ \hline
    \end{tabular}
  }
\end{table}
To evaluate the coding task-solving capabilities of our CAMEL model, specifically the LLaMA-7B fine-tuned on our comprehensive datasets, we rely on HumanEval \cite{chen2021evaluating} and HumanEval$^{+}$ \cite{evalplus}. The results, as depicted in table \ref{humaneval}, clearly demonstrate the remarkable performance of CAMEL. It surpasses not only the LLaMA-7B model but also Vicuna-7B \cite{vicuna} by a big margin. These findings underscore the critical role played by the generated datasets in enhancing LLaMA's ability to tackle coding-related tasks.


